\documentclass[a4paper]{article}
\usepackage{graphicx}
\usepackage{amsmath,amssymb}
\usepackage{hyperref}
\usepackage{minted}
\usepackage{multirow}
\usepackage{float}
\usepackage{todonotes}
\usetikzlibrary{graphs,graphdrawing}
\usegdlibrary{layered}
\usepackage{forest}
\usepackage{xstring}
\input{/home/donkarlo/Dropbox/repo/nd_ai_project/src/nd_ai/knoweldge_representation/language/formal/computer/programming/tex/latex/code_clips/external_dot.tex}
\title{}
\date{}

\begin{document}
    \maketitle
    Soar[1] is a cognitive architecture,[2] originally created by John Laird, Allen Newell, and Paul Rosenbloom at Carnegie Mellon University.
    
    The goal of the Soar project is to develop the fixed computational building blocks necessary for general intelligent agents – agents that can perform a wide range of tasks and encode, use, and learn all types of knowledge to realize the full range of cognitive capabilities found in humans, such as decision making, problem solving, planning, and natural-language understanding. \url{}
    % \bibliographystyle{plain}
    % \bibliography{/home/donkarlo/Dropbox/repo/shared_research/bibliography/bibliography.bib}
\end{document}